% Isolated post-CP26 candidate. No statement is promoted by the finite supplement.
\subsection{A first-level mixed break and the doubled tame boundary}

For the local assertions let \(p\ge5\), \(q=p^2\), \(N=q-1\).
The field \(K\) is complete discretely valued of mixed characteristic
\((0,p)\), with algebraically closed residue field, \(\mu_p\subset K\),
and absolute ramification \(e\) prime to \(p\). Put \(s=e(p+1)\),
and use the actual tame good extension \(L/K\) of degree \(N\).
Thus \(v_L(p)=eN=(p-1)s\) and \(N\mid s\).
For every two-slope group, \(B\) denotes its integral parameter
\[
 0<B<pe,\qquad \alpha=s-B,\qquad \beta=B/p,\qquad
 B\equiv c\pmod{p+1},\qquad \lambda_H=\zeta^{c-B}.
 \tag{B-RD72.1}
\]
These are parameters of actual good formal groups and actual
finite-flat quotients, with the semilinear generic descent of
Lemma~\ref{lem:B-RD69-quotients}.

\begin{lemma}[A mixed-divisibility inequality already at the first level]
\label{lem:B-RD72-first-mixed}
\status{Candidate; complete proof supplied, independent audits pending.}
Suppose two two-slope groups have isomorphic actual \(p\)-torsion,
with parameters \(C_{\mathrm{pr}}\) and \(C_{\mathrm{div}}\), where
\(p\nmid C_{\mathrm{pr}}\) and \(p\mid C_{\mathrm{div}}\).
Then necessarily
\[
                 C_{\mathrm{pr}}>C_{\mathrm{div}}.
 \tag{B-RD72.2}
\]
\end{lemma}
\begin{proof}
The first-level division field \(R/L\) of the prime-parameter group
has degree \(p\). Its high line and determinant are trivial over
\(L\), so its residual action is upper unipotent of order at most
\(p\); the nonintegral low valuation \(C_{\mathrm{pr}}/p\) makes
the action nontrivial. The actual torsion isomorphism identifies
the division fields and their unique wild fixed lines.

Fix a nonidentity \(g\in\operatorname{Gal}(R/L)\) and corresponding
low points. In the normalized valuation of \(R\), their coordinate
valuations are \(C_{\mathrm{pr}}\) and \(C_{\mathrm{div}}\), and
their nonzero displacements are high points, of valuations
\(p(s-C_{\mathrm{pr}})\) and \(p(s-C_{\mathrm{div}})\).
Formal subtraction has the same valuation as coordinate subtraction.
Writing the actual lower break as
\(r=v_R(g\varpi-\varpi)-1\), the equality and strict-inequality
parts of Lemma~\ref{lem:B-RD60-break} give
\[
 \begin{aligned}
 r&=ps-(p+1)C_{\mathrm{pr}},\\
 r&<ps-(p+1)C_{\mathrm{div}}.
 \end{aligned}
 \tag{B-RD72.3}
\]
Subtracting proves (B-RD72.2). This uses only actual first-level
torsion. No second-level comparison is imported into this lemma.
\end{proof}

\begin{proposition}[A necessary inequality for mixed second quotients]
\label{prop:B-RD72-mixed-quotients}
\status{Candidate; complete proof supplied, independent audits pending.}
Suppose an actual \(p^3\)-torsion isomorphism identifies the high
lines of two two-slope groups with
\[
 B_i=pb_i<e,\qquad qb_1<e<qb_2,\qquad p-1\nmid b_1.
 \tag{B-RD72.4}
\]
Then it is necessary that
\[
                    e(p+1)>q(pb_1+b_2).
 \tag{B-RD72.5}
\]
\end{proposition}
\begin{proof}
The actual comparison identifies the finite determinant and the
original high-line characters. Consequently
\(b_1\equiv b_2\pmod{p-1}\), as in (B-RD71.12).
Take the actual quotients \(M_i'=M_i/H_i\).
Lemma~\ref{lem:B-RD71-marked-low} gives their parameters
\(D_i=q b_i\), an actual \(p^2\)-torsion comparison, and
corresponding marked stable low lines
\[
 J_i=\operatorname{im}\bigl(M_i[p]\longrightarrow M_i'[p]\bigr).
 \tag{B-RD72.6}
\]
Their residual actions are tame, and their first-level modules are
rational over \(L\). The quotient low/high eigenvalue ratio is
\(\zeta^{(p+1)B_i}\). It is nontrivial, since
\(p-1\nmid b_i\) and \(B_i=pb_i\).
Thus each residual representation has two distinct eigenspaces.
Preservation of the marked low eigenspace forces preservation of
the other eigenspace, which is precisely its high line \(H_i'\).
This argument uses the actual marked line; it does not choose a
complement in a scalar representation.

Now quotient both \(M_i'\) by the corresponding high lines. The
inherited comparison is exactly at the first level:
\[
 (M_i'/H_i')[p]
   =\{Z\in M_i'[p^2]:pZ\in H_i'\}/H_i'.
 \tag{B-RD72.7}
\]
Since \(D_1<e<D_2\), Lemma~\ref{lem:B-RD69-quotients} gives
their respective parameters
\[
 C_1=pD_1=p^3b_1,
 \qquad C_2=s-D_2=s-qb_2.
 \tag{B-RD72.8}
\]
Both are in \((0,pe)\): the first because \(D_1<e\), and the
second because \(e<D_2<pe\). Moreover \(p\mid C_1\) and
\(p\nmid C_2\), since \(p\nmid s\).
Lemma~\ref{lem:B-RD72-first-mixed} therefore requires
\(s-qb_2>p^3b_1\). This is exactly (B-RD72.5).
All quotient identities concern actual generic torsion with its
inertia action; none asserts descent of the integral good model
through a ramified quotient.
\end{proof}

\begin{proposition}[Common-high recovery at the doubled boundary]
\label{prop:B-RD72-common-high}
\status{Candidate; complete proof supplied, independent audits pending.}
Set \(e=2N=2(p^2-1)\). An actual \(p^3\)-torsion isomorphism
between two two-slope groups which identifies their high lines
determines their finite special inertia characters.
\end{proposition}
\begin{proof}
The common determinant and high character give
\[
 c_1\equiv c_2\pmod{p-1},\qquad
 c_1-B_1\equiv c_2-B_2\pmod N.
 \tag{B-RD72.9}
\]
If either \(B_i\) is prime to \(p\),
Lemma~\ref{lem:B-RD69-rules} recovers the characters or excludes
the mixed-divisibility comparison. This includes \(B_i=e\).
Otherwise both are divisible by \(p\), so neither equals \(e\).
If both exceed \(e\), their actual high quotients have prime
parameters \(s-B_i\); their inherited \(p^2\)-comparison recovers
the characters. If they lie on opposite sides of \(e\), these
quotients have opposite parameter divisibility, excluded by the
same comparison lemma.

It remains that \(B_i=pb_i<e\). Here
\[
 1\le b_i\le2p-1,
 \qquad b_1\equiv b_2\pmod{p-1}.
 \tag{B-RD72.10}
\]
Equality of the two integers identifies \(c_1=c_2\pmod N\).
If they differ, reorder so that \(b_1<b_2\).
When \(b_1\ge2\), one has \(qb_1\ge2q>e\), and
Proposition~\ref{prop:B-RD71-separated-high} excludes this
comparison. Thus \(b_1=1\), leaving only \(b_2=p\) or
\(b_2=2p-1\).
The pair \(1,p\) exchanges the two Frobenius-conjugate special
eigenvalues by (B-RD71.14), so gives the same character.

For \((b_1,b_2)=(1,2p-1)\), all hypotheses of
Proposition~\ref{prop:B-RD72-mixed-quotients} hold:
\(p-1\nmid1\) and \(q<2(q-1)<q(2p-1)\).
The second actual quotients have
\[
 C_1=p^3,
 \qquad C_2=3p^2-2p-2,
 \qquad C_1-C_2=p(p-1)(p-2)+2>0.
 \tag{B-RD72.11}
\]
This contradicts the necessary inequality \(C_2>C_1\).
The additional arithmetic pattern is therefore excluded by an
actual first-level break comparison, completing the proof.
\end{proof}

\begin{proposition}[All whole formal comparisons at the boundary]
\label{prop:B-RD72-formal}
\status{Candidate; complete proof supplied, independent audits pending.}
At \(e=2N\), actual \(p^3\)-torsion determines the finite special
inertia character among whole potentially supersingular
height-two groups, with arbitrary finite determinant.
\end{proposition}
\begin{proof}
The one-slope and mixed-polygon cases are covered by
Proposition~\ref{prop:B-RD69-one-slope}, with no upper restriction
on \(e\). For two two-slope groups, residual wild inertia makes
their high line the unique fixed line, so the preceding
proposition applies.

For completeness consider every residual tame case at this new
endpoint. Write \(B_i=pb_i\), \(0<b_i<e\), since the first-level
modules are rational over \(L\). If both \(b_i\) are prime to
\(p\), Lemma~\ref{lem:B-RD69-rules} recovers the characters.
If exactly one is prime, say \(b_1=b\), \(b_2=pk\), choose
corresponding primitive points whose \(p\)-multiples lie outside
both high lines. Their common point field has degree \(q\)
over \(L\), forced by the first valuation \(b/q\), and all
first-level translations occur. Its normalized coordinate
valuations are \(b\) and \(pk\).
Distinct high lines would allow a translation high in the first
model and low in the second. The exact and lower-bound parts
of Lemma~\ref{lem:B-RD60-break} would give
\[
 q(\alpha_1-pk)\le b-pk.
 \tag{B-RD72.12}
\]
But \(\alpha_1=s-pb\ge e+p\), \(pk<e\), and
\(e=2q-2<p^3\). The left side is greater than \(p^3\),
whereas the right side is less than \(p^3\). Thus the high
lines correspond, and the common-high proposition applies.

Finally, if both \(b_i=pk_i\), then \(B_i=qk_i\).
Assume their high lines do not correspond. If one \(B_i>e\),
quotient that model by its high line and the other by the
corresponding stable low line. The actual \(p^2\)-comparison
has parameters \(s-B_i\) and \(B_j/p=pk_j\), of opposite
divisibility, again impossible.
Neither \(B_i=e\), since \(p\nmid e\).
If both \(B_i<e=2q-2\), they must both equal \(q\).
Lemma~\ref{lem:B-RD70-exchange} would then require
\(p-1\mid2q\), equivalently \(p-1\mid2\), impossible.
This exhausts all tame cases as well.
\end{proof}

\begin{theorem}[Exact index at twice the full tame degree]
\label{thm:B-RD72}
\status{Candidate; complete proof supplied, independent audits pending.}
For every prime \(p\ge5\), in the full B.4 monodromy setting
with height two, perfect residue field and \(\mu_p\subset K\),
\[
 \boxed{\nu_{\min}\bigl(p,2(p^2-1),2\bigr)=2.}
 \tag{B-RD72.13}
\]
In particular \(\nu_{\min}(5,48,2)=2\).
Combined with Theorem~\ref{thm:B-RD71}, the exact formula is
\[
 \nu_{\min}\bigl(p,t(p-1),2\bigr)=2
 \quad\text{for }2\le t\le2p+2,\quad p\nmid t.
 \tag{B-RD72.14}
\]
\end{theorem}
\begin{proof}
Pass to the completed maximal unramified base to prove the
upper bound. Since \(p\nmid e\), the actual cyclotomic character
modulo \(p^3\) has order \(p^2\), separating the total dimensions
\(0,1,2\) and the finite determinant as in
Theorem~\ref{thm:B-RD69}. Dimensions zero and two, and comparisons
of two systems filtered by height-one potentially good pieces,
are determined residually, with Cartier duality for dimension two.
The mixed whole-supersingular/filtered comparison is excluded by
Proposition~\ref{prop:B-RD63-mixed}, because
\(e=2q-2<p^3\). Two whole groups are handled by
Proposition~\ref{prop:B-RD72-formal}. Thus actual \(p^3\)-torsion
always determines the requested character in the full setting.
Because \(e>p\), Theorem~\ref{thm:B-RD63-arc} at \(n=2\)
supplies permitted systems with equal actual \(p^2\)-torsion and
different characters. The manuscript's least index is therefore
two. The final range adds the single endpoint \(t=2p+2\) to B71.
\end{proof}

\begin{remark}[Precise scope of the first-level obstruction]
\label{rem:B-RD72-scope}
\status{Candidate scope statement; no later-index or existence assertion.}
The new strict inequality is a first-level necessary condition,
not a first-level character-recovery theorem for arbitrary
mixed-divisibility parameters. The second quotient in
(B-RD72.7) leaves only actual \(p\)-torsion; (B-RD72.2) is
precisely the assertion used there.
The intermediate residual eigenvalues are distinct in the new
boundary pair, so their preserved low lines determine the high
lines. Scalar intermediate actions are not treated by that
argument; the separate scalar-safe B71 proposition is used in
the cases where it applies.
No absolute ramification divisible by \(p\) is settled here.
The next prime-to-\(5\) index \(e=52\), and the wild index
\(e=20\), remain outside this theorem.
The retained exact supplement
\path{code/reduction_descent/verify_first_level_mixed_boundary.py}
reconstructs both actual high-quotient Newton polygons from the
universal vertices and counts quotient roots using the product map.
Across twelve primes from \(5\) through \(43\), it checks
\(24\) double-quotient root records, \(264\) corresponding
distinct-eigenvalue records, \(545{,}376\) first-level prime breaks,
and \(1{,}074{,}792\) tame mixed point-field inequalities.
It also checks the entire stated small common-norm case split,
the boundary strict-break contradiction at every tested prime,
and all remaining tame quotient-divisibility branches.
For \(p=5,e=48\), the actual second parameters are \(125\) and
\(63\); the prime model requires lower break \(1062\), whereas
the divisible model requires it to be strictly less than \(690\).
The author run passed in \(0.5185015480092261\) seconds at
22:14:11 UTC. The complete program, output, and pre-certificate
proof are retained. No executed numerical check failed.
These finite checks supplement the all-prime proof and do not
provide a formal-group existence test.
\end{remark}
